% Generated draft: requires mathematical and source review.
\documentclass[11pt]{article}
\usepackage[margin=25mm]{geometry}
\usepackage{fontspec}
\setmainfont{Latin Modern Roman}
\usepackage{amsmath,amssymb,mathtools,unicode-math}
\setmathfont{Latin Modern Math}
\usepackage{enumitem,needspace,xurl,hyperref}
\setlength{\parindent}{0pt}
\setlength{\parskip}{5pt}
\setlength{\emergencystretch}{3em}
\newcommand{\N}{\mathbb{N}}
\newcommand{\Z}{\mathbb{Z}}
\newcommand{\Q}{\mathbb{Q}}
\newcommand{\R}{\mathbb{R}}
\newcommand{\C}{\mathbb{C}}
\newcommand{\F}{\mathbb{F}}
\newcommand{\A}{\mathbb{A}}
\newcommand{\PP}{\mathbb{P}}
\newcommand{\E}{\mathbb{E}}
\newcommand{\eps}{\varepsilon}
\newcommand{\dd}{\,\mathrm{d}}
\begin{document}
\section*{UnsolvedMath ID: 20002702}
\textbf{Source problem code:} catalogue entry
\subsection{Definitions and mathematical statement}
\textbf{Source-fragment identity and underlying question.} The aid identifies this record as a page-broken continuation of the final bullet of Problem 1.2, not a separate Problem 1.1: the apparent number belongs to a parenthetical cross-reference. Permanent ID and source code are unchanged. The fragment asks how nonnesting partitions and root order extend to noncrystallographic finite Coxeter types, what essential barriers arise for familiar Lie-algebra and affine-hyperplane constructions, and what objects Chapoton's proposed \(H\)-triangle counts. Its preceding antecedent ``these objects'' is not fully supplied. The text mentions Kriloff--Ram [33] and Chapoton [18], without the bibliography entries.

In crystallographic root systems, a standard positive-root order has \(\alpha\leq\beta\) when \(\beta-\alpha\) is a nonnegative integer combination of simple roots, and nonnesting objects are antichains. Extending this formula or its combinatorial consequences to irrational root coefficients is part of the problem, not a fixed definition in the source. The source's suggested barriers are research questions, not an impossibility theorem for every generalized Lie or affine construction.

\textbf{Reported uniform dihedral model.} For each integer \(m\geq3\), define an abstract poset with \(m\) positive-root labels: two incomparable simple elements \(\alpha_1,\alpha_2\), both below \(\beta_1\), and a chain
\[
 \beta_1<\beta_2<\cdots<\beta_{m-2}.
\]
This fixes the V-then-chain poset up to isomorphism; it does not claim uniqueness among all possible noncrystallographic root-poset axiom systems. With the convention
\[
 H(s,t)=\sum_{A\text{ antichain}}s^{|A\cap\{\alpha_1,\alpha_2\}|}t^{|A|},
\]
the aid's exact enumerator is
\[
 H(s,t)=1+2st+(m-2)t+s^2t^2.
\]
The empty antichain, simple singleton antichains, nonsimple singleton antichains and the two-simple antichain account for all terms.

The rank-two \(I_2(m)\) cluster complex is represented by a cycle on \(m+2\) vertices, with two adjacent distinguished negative vertices and \(m\) positive vertices. Weight a face by \(x^{\#\mathrm{positive}}y^{\#\mathrm{negative}}\). The independently counted face polynomial is
\[
 F(x,y)=1+mx+2y+(m-1)x^2+2xy+y^2.
\]
Its edges comprise \(m-1\) positive-positive, two mixed and one negative-negative edge. The aid reports the exact convention-sensitive Chapoton relation
\[
 H(s,t)=(1-t)^2F\left(\frac{t}{1-t},\frac{st}{1-t}\right),
\]
\[
 F(x,y)=(1+x)^2H\left(\frac yx,\frac{x}{1+x}\right).
\]
These are polynomial identities after cancellation; the displayed inverse rational substitution is first evaluated where defined. The \(H\) variables count simple roots and total size in the stated order, which is the swap needed relative to the alternative standard convention. This works uniformly across crystallographic and noncrystallographic \(I_2(m)\), independently of parity or root-length ratio; it is an abstract enumerative model rather than an identification with every geometric order.

The aid additionally records a positive \(H_3\) construction and Cuntz--Stump and Chen--Kriloff \(H_4\) obstructions, but gives neither complete references nor the exact root-poset axioms/requirements involved. Those obstructions are not promoted to nonexistence of every possible \(H_4\) counting object. The dihedral identity does not resolve the full extension program or the essential-barrier question in all noncrystallographic types.
\subsection{Short English statement}
This is a split continuation of a broader Coxeter-theory problem. The aid gives a complete dihedral antichain/cluster-cycle identity with explicit variable conventions, while preserving the incomplete \(H_3\)/\(H_4\) literature qualifications and the unanswered generalization/barrier questions.
\subsection{Sources}
\begin{itemize}
\item Local UnsolvedMath registry, permanent ID \(20002702\), catalogue entry, \texttt{aim-probability-notes.json}, \texttt{notes[143]}, workshop \emph{Braid groups, clusters and free probability}: \url{https://aimath.org/WWN/braidgroups/braidgroups.pdf}. The aid identifies the fragment as the continuation of Problem 1.2; the source item field reads 1.1.
\item Original attributions: Cathy Kriloff and Arun Ram [33] on representation-theoretic noncrystallographic issues; Fr\'ed\'eric Chapoton [18] on the \(H\)-triangle. Bibliography entries are not supplied.
\item Supplied research summary records an \(H_3\) construction and Cuntz--Stump and Chen--Kriloff \(H_4\) obstructions without full references/axioms, and supplies the dihedral coefficient identity. Attempt 1: context\_only, not\_a\_problem, rigor ``proved'', low novelty confidence; registry open/OPEN-TRIAGE. These source-fragment labels do not erase the underlying questions. The machine-generated aid is not peer reviewed and requires independent verification before citation; no PDF visit is claimed.
\end{itemize}
\end{document}
